\section{Introduction}\label{sec:introduction}

Independent diagonal disorder changes both the spectrum of a graph and
the spatial distribution of its eigenvectors.  Anderson introduced this
model to describe how randomness in a medium can prevent electronic
transport~\cite{Anderson1958}.  On the infinite regular tree, or Bethe
lattice, removing an edge separates independent branches.  The resulting
recursive distributional equations for Green functions underlie the
self-consistent theory of Abou-Chacra, Thouless, and
Anderson~\cite{AAT1973} and make the tree a basic setting for studying
the transition between localized and extended states.

The weak-disorder regime on the tree has a rigorous foundation.
Klein~\cite{Klein1994,Klein1998} proved purely absolutely continuous
spectrum on compact subintervals of the clean spectral band for
sufficiently small disorder, for every regular-tree degree at least
three.  He also proved that the integrated density of states is
continuously differentiable there, so the density of states is
continuous.  Aizenman, Sims, and Warzel~\cite{ASW2006} developed a
fluctuation approach to the stability of absolutely continuous spectral
weight as disorder tends to zero.  Froese, Hasler, and
Spitzer~\cite{FHS2007} gave a geometric proof for the degree-three tree,
using the hyperbolic geometry of the Green-function recursion.  We
establish within the paper the quantitative recursion stability and
boundary regularity required for the finite-graph argument.

A random regular graph of fixed degree is locally a tree, while its
finite spectrum raises a further question: do eigenvalues in a window
of length comparable to the inverse graph size have the statistics of
a real symmetric Gaussian matrix?  The present paper answers this
question in the clean bulk for sufficiently weak positive disorder that
remains fixed as the graph grows.  The conclusion holds at each fixed
energy and identifies the entire microscopic point-process law.

For the clean model, the tree spectral measure and the limiting empirical
density are the Kesten--McKay law~\cite{Kesten1959,McKay1981}.
Local laws and eigenvector delocalization were established for growing
degrees by Bauerschmidt, Knowles, and Yau~\cite{BKY2017Local}, and for
sufficiently large fixed degrees by Bauerschmidt, Huang, and
Yau~\cite{BHY2019Local}.  Huang and Yau~\cite{HY2024Spectrum} extended
the fixed-degree theory to every degree at least three, including the
spectral edges.  Microscopic statistics require further input:
Bauerschmidt, Huang, Knowles, and Yau~\cite{BHKY2017} proved bulk GOE
gap statistics and energy-averaged correlation statistics for
a range of polynomially growing degrees.  Huang, McKenzie, and
Yau~\cite{HMY2024Edge} established optimal rigidity and edge universality
at every fixed degree at least three.  The companion~\cite{bu} proves
fixed-energy bulk universality for the clean fixed-degree model.  We
adapt its combinatorial entropy and Gaussian insertion arguments below.

For positive disorder, Anantharaman and
Sabri~\cite{AnantharamanSabri2017} proved quantum ergodicity for the
Anderson model on expanding regular graphs converging locally to the
regular tree, at weak disorder in compact subsets of the clean band.
Their general framework~\cite{AnantharamanSabri2019} relates spatial
equidistribution to quantitative regularity of tree Green functions.
These conclusions concern averaged spatial equidistribution.  The
present theorem concerns joint eigenvalue statistics in a microscopic
window around a fixed energy.

The eigenvector part of our proof follows the entropy method of
Backhausz and Szegedy~\cite{BackhauszSzegedy2019}.  They connect almost
eigenvectors of finite random regular graphs to Gaussian waves on the
tree through local type counts, Gaussian smoothing, weighted
Fisher-information inequalities, and their equality cases.  Their
limiting Gaussian law allows loss of variance under weak convergence.
Here we separately preserve second moments in the microscopic limit
and adapt the entropy argument to laws conditioned on the random
potential.  The conditional covariance is determined by the random tree Green
function.

The final identification uses the dynamics introduced by
Dyson~\cite{Dyson1962}.  The local relaxation method of Erd\H{o}s,
Schlein, and Yau~\cite{ESY2011} and the homogenization argument of
Bourgade, Erd\H{o}s, Yau, and Yin~\cite{BEYY2016} established key links
between short-time Dyson Brownian motion and universal local statistics.
For deterministic initial data with a regular local density, Landon and
Yau~\cite{LandonYau2017} proved short-time gap and energy-averaged
universality, and Landon, Sosoe, and Yau~\cite{LSY2019} proved the
fixed-energy version.  We apply the latter theorem after comparing
the effect of measure-preserving potential changes with a Gaussian
perturbation and verifying the required regularity and density
normalization.

\subsection{Model and main result}

Fix an integer $d\ge3$.  For $n>d$ with $nd$ even, let
$\mathcal G_{n,d}$ be uniform among all simple labelled $d$-regular
graphs on $\{1,\ldots,n\}$, and let $A_{n,d}$ be its adjacency matrix.
Let $(\omega_v)_{v=1}^n$ be independent uniform random variables on
$[-1,1]$, independent of the graph.  For $w>0$, put
\[
 H_{n,d,w}=A_{n,d}+w\diag(\omega_1,\ldots,\omega_n).
\]
We use ordinary adjacency, without dividing by $\sqrt{d-1}$, and count
all eigenvalues $\lambda_1,\ldots,\lambda_n$ with multiplicity.
The graph law has no further conditioning.

Let $\mathbb T_d$ be the infinite $d$-regular tree with root $o$, equipped
with independent uniform potentials $\omega_v\in[-1,1]$.  Its density-of-states
measure is the deterministic probability measure
\begin{equation}\label{eq:tree-dos-definition}
 \nu_{d,w}(B)=
 \EE\left\langle\delta_o,
 \mathbf1_B\bigl(A_{\mathbb T_d}+w\diag(\omega_v)\bigr)\delta_o
 \right\rangle,
 \qquad B\subset\RR\ \text{Borel}.
\end{equation}
This definition fixes the unfolding density directly from the tree
operator.  Write $M_n$ for normalized GOE: its entries on and above
the diagonal are independent centered real Gaussians, with variance
$1/n$ off the diagonal and $2/n$ on the diagonal.  If its eigenvalues
are $\mu_1,\ldots,\mu_n$, the unit-intensity bulk reference at zero is
\[
 \Xi_n^{\GOE}=\sum_{i=1}^n\delta_{(n/\pi)\mu_i}.
\]

\begin{theorem}[Fixed-energy universality]\label{thm:main}
For every $d\ge3$ and every $0<\kappa<2\sqrt{d-1}$, there is
$w_0=w_0(d,\kappa)>0$ with the following property.  For each fixed
$0<w<w_0$, the restriction of $\nu_{d,w}$ to an open neighborhood of
\[
 I_{d,\kappa}=[-2\sqrt{d-1}+\kappa,\,2\sqrt{d-1}-\kappa]
\]
has a continuous density $\rho_{d,w}$, strictly positive on
$I_{d,\kappa}$.  For every fixed $E\in I_{d,\kappa}$, set
\[
 \Xi_{n,E}=\sum_{i=1}^n
       \delta_{n\rho_{d,w}(E)(\lambda_i-E)}.
\]
Then, for every nonnegative $f\in C_c(\RR)$,
\begin{equation}\label{eq:main-laplace}
 \left|
 \EE\exp\left(-\int f\,d\Xi_{n,E}\right)
 -\EE\exp\left(-\int f\,d\Xi_n^{\GOE}\right)
 \right|\longrightarrow0
\end{equation}
as $n\to\infty$ through all sizes $n>d$ with $nd$ even.
\end{theorem}

The disorder strength, degree, energy, and test function are fixed before
the graph size tends to infinity.  The conclusion identifies the full
microscopic point-process law.  The tree density in
\eqref{eq:tree-dos-definition} also determines the variance profile of
the eigenvector marks used in the proof.

\subsection{Proof strategy and the role of the environment}

The first task is to compare finite-graph resolvents with random tree
resolvents.  A branch Green function satisfies a scalar recursion whose
inputs come from independent child branches.  Section~\ref{sec:tree}
proves stability of this recursion near the clean bulk solution, including
control of values close to the boundary of the upper half-plane.
Section~\ref{sec:local-law} then derives an approximate recursion for
the empirical law of finite-graph cavity resolvents.  It uses switches
of a bounded number of edges, spectral averaging in the potentials, and
the Ward identity.  The result is a distributional local law: the random
environment remains visible in the limiting resolvent.

The second task is to identify eigenvector marks at sampled vertices.
Normalize an eigenvector $u_i$ by $\|u_i\|^2=n$ and consider a fixed
finite list of eigenvalues within distance $O(n^{-1})$ of $E$.
In a subsequential local limit, the corresponding coordinates at a
sampled root satisfy the tree eigen-equation.  Microscopic uniform
integrability preserves their exact second moments.  A small simultaneous
change of all potentials, combined with spectral averaging, eliminates
conditional means.  Section~\ref{sec:entropy} establishes an entropy
inequality for these marks conditional on the entire tree environment.

The Gaussian comparison field has covariance given by the imaginary part
of the tree Green function at $E+i0$.  Like the eigenvector marks, it
satisfies the tree eigen-equation, which imposes one linear constraint
for each eigenvector coordinate on a star.  On that constraint space,
its precision splits into contributions
from the incident edges; the two contributions at either end of an edge
add to the edge precision.  Section~\ref{sec:rigidity} uses this identity
and re-rooting invariance in a Fisher-information comparison.  After
removing the directions normal to the constraint space and subtracting
the Gaussian entropy constants, the counting inequality makes the
entropy difference
nonnegative, while the comparison makes it increase to zero as Gaussian
noise is added.  Equality is therefore forced, and its equality case gives
independent Gaussian coordinates at the root, conditional on the
environment, with variance
\[
 \frac{\Im m_o(E+i0)}{\pi\rho_{d,w}(E)},
 \qquad
 m_o(\zeta)=\left\langle\delta_o,
 (A_{\mathbb T_d}+w\diag(\omega_v)-\zeta)^{-1}\delta_o\right\rangle.
\]
This conditional statement retains information that would be lost by
replacing the random variance by its average.

The final task is to determine the microscopic eigenvalue process.
Section~\ref{sec:resolvents} represents sampled resolvents by Gaussian
marks on that process, with the real additive constants fixed by the
joint local law.  Resampling one potential gives an exact rank-one
identity.  Together with Gaussian concentration, it improves the initial
slow estimates to uniform regularity of the limiting Stieltjes transform.
Measure-preserving exchanges of short intervals in the potential law
then provide many small perturbations.  Their number increases inside the
already constructed limiting law.  After truncating distant microscopic
positions, the switch identity represents their effect by a sum of small
signed rank-one matrices formed from the Gaussian marks.
Section~\ref{sec:insertion} replaces the centered sum by GOE noise.
Because the potential exchanges preserve the ensemble, this construction
approximates each subsequential point-process law by the eigenvalues of a
finite diagonal matrix of the retained positions, plus a scalar shift and
GOE noise.  This comparison is made after averaging over the sampled
environments.  Section~\ref{sec:dbm} then applies fixed-energy Dyson
Brownian motion to these matrices, checking their regularity and density
normalization.  The resulting unfolding is exactly the one in
Theorem~\ref{thm:main}.

The distributional local law, the environment-conditioned entropy
argument, and the use of potential resampling to improve microscopic
regularity are the principal intermediate constructions.  Their proofs
keep the random tree response explicit until the final Gaussian
comparison.

\subsection{Conventions}

We write $V_v=w\omega_v$, so that potentials denoted by $V$ include
the coupling constant and are uniform on $[-w,w]$.  We suppress $n,d,w$
in $H$ when they are fixed, and use
$G(\zeta)=(H-\zeta)^{-1}$ for $\zeta\in\HH$; imaginary parts of diagonal
resolvents are positive.  Set $b=d-1$.  Choose once and for all a compact
interval $J$ in $(-2\sqrt b,2\sqrt b)$ whose interior contains
$I_{d,\kappa}$.  All small-coupling choices will be uniform for energies
in $J$.  Constants in finite-graph estimates may depend on the fixed
positive $w$.

The notation $W_2$ denotes quadratic Wasserstein distance for the specified
Euclidean metric.  Conditional laws of colors in the entropy argument
refer to the finitely many eigenvector coordinates; these colors are
independent Gaussian coordinates only after that conclusion is proved.
An environment consists of the potentials, including their coupling
constant.  Probability and expectation also include auxiliary uniform
vertex or edge samples when these have been introduced.
